Going back to the completed-versus-abandoned split from Section V-A: LTR only looks fast on average because it finishes the easy requests and lets the hard ones time out. Our scheduler finishes far more requests overall, including many of the big ones that LTR would drop. That's exactly why our average ends up higher, since it isn't doing a worse job but just finishing harder work. So can a scheduler be fast and fair at the same time, even with bad predictions? Not fully, and not on one single number. But the fairness cost stays small, and it shrinks even more as the server gets busier. That matches with what we saw across the four scenarios in Section V, where our advantage over FCFS kept growing from Light all the way to Saturation. Here's why: under light load, almost every policy finishes quickly no matter how it orders the queue, so aging doesn't have much to prove. Under heavy load, the order actually matters, and that's where aging earns its keep for the system performance.

This same pattern, which gets worse at higher load, shows up again when we compare our results with prior work. We agree with \cite{kumar2024latency} in one point: FCFS comes out worst at every load level and with every kind of noise we tested. Where we differ is on LTR. In our original, simple noise model, LTR never actually gets worse than FCFS, which doesn't match with the sharper failure \cite{kumar2024latency} describes for real predictors facing prompts they weren't trained on. Part of the reason is how our noise is built. Even at 80 percent error, the ranking that our noisy predictions produce still tracks the true request sizes closely, so a genuinely large request rarely gets mistaken for a small one. LTR keeps making mostly sensible calls even when it's badly informed. That's why we built a second, harder noise model for Section V-E, where some requests get a prediction that has almost nothing to do with their real size, instead of just a noisy one. Once we ran that version, LTR's completion time went up by 119 percent. That's much closer to what \cite{kumar2024latency} and \cite{fu2024ltr} both report, and it suggests that the kind of failure that actually hurts LTR isn't random error. It's more like a predictor running into something it was never trained on.

We want to be upfront about three limitations. First, our simulator only runs one request per time, so there's no real concurrent batching. That was a scope decision from the start, not something we missed, because adding batching would mean modeling GPU-level effects we chose not to touch here, and it follows the same simulation-first approach used in \cite{agrawal2024vidur}. Second, we didn't build a real memory pool. The preemption count in our code only tracks how often a predicted size turned out too small at admission time; it doesn't represent an actual memory eviction \cite{kwon2023pagedattention}. Third, we don't model prefill and decode as separate phases the way \cite{sheng2024fairness} and \cite{patel2024splitwise} do, and we only look at a single server, not several working together like in \cite{sun2024llumnix}.

These same gaps point pretty directly to where the next six months should go. The first step is giving to the engine a real memory pool with a fixed size, so preemption stops being just a warning sign and becomes a real event, where a request gets kicked out or swapped once memory runs out. The second is making alpha and beta adjust by themselves as load changes, instead of picking one fixed setting per scenario: raise beta when the queue builds up, lower it when things are calm. This follows directly from what we found in Section V-B, where each scenario needed a different beta to reach the same fairness target. A scheduler that senses its own load and adjusts by itself could probably reach that target without needing to be re-tuned by hand every time things change.

Past these two, swapping the synthetic noise generator for a real trained predictor, closer to what \cite{kumar2024latency} and \cite{fu2024ltr} used, would tell us if the formula still works well against real prediction errors and not just simulated ones. There's also a related but different question: what happens if the noise is not honest at all? If a user finds out how the formula scores requests, they could start shaping their prompts to fake a higher priority. None of our current noise models, symmetric or out-of-distribution, actually catches that. Running our numbers against an outside simulator like \cite{cho2024llmservingsim} or \cite{agrawal2024vidur} would also help to confirm our results are not just an effect from our own code. One direction worth to name, even if we don't pursue it, is letting a language model act as the scheduler itself, as in \cite{abgaryan2024llmschedule}. That would be a real contrast with the fixed formula we built here.